\documentclass{ximera}

% MATH -----------------------------------------------------------
\newcommand{\nats}{\mathbb N}
\newcommand{\ints}{\mathbb Z}
\newcommand{\rats}{\mathbb Q}
\newcommand{\reals}{\mathbb R}
\newcommand{\complex}{\mathbb C}
\newcommand{\powerset}{\mathscr P}

%-------------------------------------------------------------

\pgfplotsset{compat=1.18}
\begin{document}
\begin{abstract}
 \end{abstract}

    \title{Exemplar Examples 3}
    
    \maketitle


Let's consider the differential equation $y^{\,\prime}=y^{2/3}e^{x/3}$.\\


\begin{enumerate}

\item Does this equation have any constant solutions?\\ \\

% YES!  y=0

\item Let's write this equation the form $h(y)y^{\,\prime}=g(x)$.

\vfill

% y^(-2/3) y^{\,\prime}=e^(x/3) when y \neq 0

\item Let's integrate both sides with respect to $x$.

% 3y^(1/3)=3e^(x/3)+C  -

\vfill


\item Now let's solve this differential equation for its general solution (including any constant solutions).

% y=(e^(x/3)+C)^3 or y=0


\vfill

\vfill

\vfill

\item  From the above, let's determine TWO solutions to the IVP $y^{\,\prime}=y^{2/3}e^{x/3}$ $y(0)=0$.

% y=(e^(x/3)-1)^3 and y=0.

% SANITY CHECK: y^{\,\prime}=e^(x/3)(e^(x/3)-1)^2 =y^(2/3)e^(x/3) checks out!

\vfill

\newpage

\item Why doesn't this contradict the Existence and Uniqueness Theorem?

\end{enumerate}



\end{document}